% Generated draft: requires mathematical and source review.
\documentclass[11pt]{article}
\usepackage[margin=25mm]{geometry}
\usepackage{fontspec}
\setmainfont{Latin Modern Roman}
\usepackage{amsmath,amssymb,mathtools,unicode-math}
\setmathfont{Latin Modern Math}
\usepackage{enumitem,needspace,xurl,hyperref}
\setlength{\parindent}{0pt}
\setlength{\parskip}{5pt}
\setlength{\emergencystretch}{3em}
\newcommand{\N}{\mathbb{N}}
\newcommand{\Z}{\mathbb{Z}}
\newcommand{\Q}{\mathbb{Q}}
\newcommand{\R}{\mathbb{R}}
\newcommand{\C}{\mathbb{C}}
\newcommand{\F}{\mathbb{F}}
\newcommand{\A}{\mathbb{A}}
\newcommand{\PP}{\mathbb{P}}
\newcommand{\E}{\mathbb{E}}
\newcommand{\eps}{\varepsilon}
\newcommand{\dd}{\,\mathrm{d}}
\begin{document}
\section*{UnsolvedMath ID: 30004472}
\textbf{Source problem code:} catalogue entry
\subsection{Definitions and mathematical statement}
Let \(X\) be a smooth projective variety of pure dimension \(d\), equipped with an ample line bundle \(L\) defining its polarization, and let \(\mathcal F\) be an \(\ell\)-adic local system on \(X\). The supplied record does not specify the base field, the coefficient field, whether cohomology is geometric, or the meaning of ``relevant'' in its original restriction to relevant local systems. Consequently, the following is a formulation schema whose intended hypotheses require clarification.

For the usual étale interpretation, write \(E\) for the chosen \(\ell\)-adic coefficient field, require \(\ell\) to be invertible in the base field, and interpret a local system as a finite-rank lisse \(E\)-sheaf. These conventions specify the notation; they do not resolve the source's missing hypotheses. Let
\[
\eta=c_1(L)\in H^2_{\mathrm{\acute et}}(X,E(1)),
\qquad
\mathcal F(i)=\mathcal F\otimes_E E(i),
\]
where \(E(i)\) denotes the \(i\)-th Tate twist. If geometric cohomology is intended, \(X\) in these cohomology groups must instead denote the base change to an algebraic closure of its base field.

The supplied local-system question asks whether, for every polarized smooth projective \(d\)-fold \(X\), every relevant \(\ell\)-adic local system \(\mathcal F\), and every integer \(i\) with \(0\leq i\leq d\), cup product induces an isomorphism
\[
\eta^i\smile-:\ H^{d-i}_{\mathrm{\acute et}}(X,\mathcal F)
\longrightarrow
H^{d+i}_{\mathrm{\acute et}}(X,\mathcal F(i)).
\]
The range \(0\leq i\leq d\) makes explicit the conventional hard Lefschetz range; it is not stated in the supplied formulas. The supplied formulas also omit the Tate twist. Displaying the untwisted target \(H^{d+i}(X,\mathcal F)\) requires an explicitly justified convention or an identification of twists, which the record does not provide.

No semisimplicity, arithmeticity, or rank restriction is imposed in the clean statement. Such restrictions must not be silently added. The original wording also asks about ``all \(\ell\)-adic sheaves,'' whereas its displayed question concerns local systems; the accompanying literature assessment instead discusses arbitrary semisimple perverse \(\ell\)-adic sheaves. These are different classes of coefficients and do not determine a single precise conjecture without further source clarification. The displayed schema records the local-system target and leaves these qualifications unresolved.
\subsection{Short English statement}
The supplied question asks whether powers of a polarization give hard Lefschetz isomorphisms for every relevant \(\ell\)-adic local system on a smooth projective variety. The base field, coefficients, geometric cohomology convention, and permitted local systems remain unspecified. Its references to all sheaves and to semisimple perverse sheaves require clarification. The supplied open label is dated curation, not an independently verified conclusion.
\subsection{Sources}
\begin{itemize}
\item Local registry: \texttt{data/unsolvedmath/problems.json}, supplied record with permanent UnsolvedMath ID \(30004472\), source problem code \texttt{catalogue entry}, in the Oberwolfach Reports Open Problems collection. No individual proposer is supplied; the proposed year is recorded as \(2020\).
\item Supplied report citation: \emph{Arithmetic Geometry (2021)}, Oberwolfach Reports, DOI \texttt{10.4171/owr/2020/20}, pp.\ 1023--1082. Original OWR report: \url{https://doi.org/10.4171/owr/2020/20}. The background describes these reports as documenting mathematical workshops and their open-problem discussions. The cited report has not been independently consulted for this draft.
\item Supplied literature reference labeled ``Esnault--Kerz'': \url{https://doi.org/10.1007/s00029-021-00681-y}.
\item Supplied literature reference labeled ``rank-one hard Lefschetz'': \url{https://arxiv.org/abs/1908.08291}.
\item Literature assessment dated exactly \texttt{2026-08-22}: the supplied curation labels the problem open and states that rank-one and arithmetic semisimple cases of hard Lefschetz are proved, while the conjectural statement for arbitrary semisimple perverse \(\ell\)-adic sheaves, and hence the broader all-sheaves question, remains unresolved. This is supplied dated curation, not independently verified here. The background explicitly requests independent verification of time-sensitive status and cited claims.
\item The registry labels the statement \texttt{corrected\_verified} and reports a manual semantic, typographic, and LaTeX review claiming that the clean statement is self-contained and preserves the target. These are supplied editorial assertions; the formulation issues recorded in this draft remain unresolved.
\end{itemize}
\end{document}
